\par\noindent\textbf{(1)} 构造三态复乘积 \(\mathcal B=\langle\psi_1|\psi_2\rangle\langle\psi_2|\psi_3\rangle\langle\psi_3|\psi_1\rangle\)。每个内积包含振幅与相位，定义闭合Pancharatnam相位为 \(\gamma=\arg\mathcal B\)，沿短大圆弧连接三个偏振态。
\par\noindent\textbf{(2)} 各Jones向量可以独立改相位 \(|\psi_j\rangle\mapsto e^{i\alpha_j}|\psi_j\rangle\)。三内积分别多出 \(e^{i(\alpha_2-\alpha_1)}\)、\(e^{i(\alpha_3-\alpha_2)}\)、\(e^{i(\alpha_1-\alpha_3)}\)，乘积为1，故 \(\mathcal B\) 的相位严格不受局部规范选择影响。
\par\noindent\textbf{(3)} 纯偏振态在Poincaré球上表示为Bloch单位矢量。对相邻态按最短测地线平行输运的闭合三角形，几何相位为 \(\gamma=-\Omega_s/2\)（具体正负取决于Jones/Stokes球取向约定），\(\Omega_s\) 是有向球面固角。其几何意义与自旋 \(1/2\) Berry相位一致。
\par\noindent\textbf{(4)} 若任意相邻态正交，则相应内积为零，\(\mathcal B=0\)，其 \(\arg\) 不定义；靠近这种点相位可能剧变。物理上完全正交偏振场不形成通常意义的相干干涉相位，不能给零幅复数指派任意连续相位。
